\documentstyle{article}
\begin{document}
%\bibliographystyle{/u/caphorn/0/saga/merlet/LaTeX/englishplain}

\begin{center}
Forward kinematics of parallel robots\\
J-P. Merlet\\
INRIA Sophia-Antipolis, France
\end{center}

\vspace{1cm}

{\bf Abstract:} The purpose of this paper is to present a state of the
art in the field of the forward kinematics of parallel robots.
The problem is to determine the possible pose of the end-effector of
the robot being given its articular coordinates and amount to solve
the non-linear system of equations obtained from the inverse
kinematics. This problem is known to be one of the most challenging in
the field of kinematics: it involves the resolution of a
relatively large set of non-linear equations in order to determine the
current pose of the end-effector. Although they are many different
possible mechanical architecture for parallel robots the direct
kinematics of most of them may be reduced to solving the direct
kinematics for a small subset of architectures, the most important
being the famous Gough platform.

This problem may be considered from different viewpoints:
\begin{enumerate}
\item determining the maximum number of solutions of the system either
real or complex
\item determining the maximum number of real solutions of the system
\item calculating all the solutions without constraint on the
computation time
\item finding an example which has exactly the maximum number of real
solutions calculated for problem 2
\item if all the real solutions have been found, sort the solutions
set in order to determine the current pose
\item calculating the current pose with real-time constraint
\end{enumerate}

The two first problems have been addressed in the last ten years,
starting with particular geometries of robots (see for
example~\cite{charentus89,griffis89,hunt78,nanua89,faugere95})
until a solutions
has been given for the most general cases of Gough platform
(see~\cite{lazard92,mourrain93,raghavan91,ronga92}, for which they may
be up to 40 solutions.

Problem 3 has received a lot of attention and many algorithms have
been proposed for particular architectures. Still for the most general
case of Gough platform although specific algorithms have been
proposed~\cite{husty94-1} they are not exactly satisfactory either
because their computation time is still large or they have to rely on
symbolic computation system to find the solutions.

Problem 4 is still an unsolved problem although a conjecture is that
there will always be a robot geometry and a set of articular
coordinates such that the number of real solutions of the direct kinematics
will reach the maximum. This has been shown for particular
architectures~\cite{lazard94,merlet89-1} and recently for the most
general case of general Gough platform~\cite{dietmaier98}. Still for
many architectures the conjecture has still to be
verified~\cite{faugere95}.

Problem 5 is quite important in practice: for real applications
solving the direct kinematics means to determine the current, real
pose of the end-effector of the robot. Consequently computing all the
solutions is only the first step of the resolution of the direct
kinematics. Then is it is necessary to sort the solutions in order to
determine the current pose of the end-effector. Currently there is no
known solution to this problem: sorting the solution will have to rely
on an algorithm which will be able to determine which solution may be
reached from an initial assembly mode without crossing a singularity
(but this is not sufficient as different solutions may be joined
without crossing a singularity as shown
in~\cite{innocenti92-3,hunt98}) and without having to dismantle the
robot.

Problem 6 is also very important in practice. Usually a Newton-Raphson
scheme~\cite{reboulet85}
 is used to solve this problem, as an estimate of the
solution is known, but this method has some drawbacks. The first
problem in this
method is to determine the parameterization of the end-effector pose
(especially the orientation part) in order to have the largest
convergence domain. The second drawback is that Newton method does not
always converge toward the solution which is the closest to the
estimate: this is a major drawback as the control of the robot heavily
rely on the direct kinematics. Alternative method have been proposed:
neural network~\cite{geng91,yee91}, genetic
algorithms~\cite{boudreau96}, interval analysis~\cite{didrit98},
improved
Newton method~\cite{chetelat97} and
their merits will be discussed.

%\bibliography{/u/caphorn/0/saga/merlet/Texte/robot_parallele,/u/caphorn/0/saga/merlet/Texte/mecanisme}
\newcommand{\vide}{\setcounter{DATE}{-1}} \newcommand{\private}{Communication
  personelle}
\newcounter{DATE}
\newcommand{\Janvier}{\setcounter{DATE}{1}}
\newcommand{\Fevrier}{\setcounter{DATE}{2}}
\newcommand{\Mars}{\setcounter{DATE}{3}}
\newcommand{\Avril}{\setcounter{DATE}{4}}
\newcommand{\Mai}{\setcounter{DATE}{5}}
\newcommand{\Juin}{\setcounter{DATE}{6}}
\newcommand{\Juillet}{\setcounter{DATE}{7}}
\newcommand{\Aout}{\setcounter{DATE}{8}}
\newcommand{\Septembre}{\setcounter{DATE}{9}}
\newcommand{\Octobre}{\setcounter{DATE}{10}}
\newcommand{\Novembre}{\setcounter{DATE}{11}}
\newcommand{\Decembre}{\setcounter{DATE}{12}}
\newcommand{\summer}{\setcounter{DATE}{13}}
\newcommand{\frenchmonth}{
\ifcase\value{DATE}\or January \or February \or March \or
April \or May \or June \or July \or August \or September \or October \or
November \or December \or Summer \fi}
\newcommand{\frenchmonthh}{
\ifcase\value{DATE}\or January, \or February, \or March, \or
April, \or May, \or June, \or July, \or August, \or September, \or October,
\or
November, \or December, \or Summer \fi}
\newcommand{\frenchdate}[3]{#2\frenchmonthh #1, \space #3}
\newcommand{\frenchmy}[2]{#1\frenchmonth \space #2}
\newcommand{\frenchddate}[5]{#3\frenchmonth #1#4-\frenchmonth #2, \space #5}
\begin{thebibliography}{10}

\bibitem{boudreau96}
Boudreau R. and Turkkan N.
\newblock Solving the forward kinematics of parallel manipulators with a
  genetic algorithm.
\newblock {\em J. of Robotic Systems}, 13(2):111--125,
  \frenchmy{\Fevrier}{1996}.

\bibitem{charentus89}
Charentus S. and Renaud M.
\newblock Calcul du mod\`ele g\'eom\'etrique direct de la plate-forme de
  {Stewart}.
\newblock Research Report 89260, LAAS, Toulouse, France,
  \frenchmy{\Juillet}{1989}.

\bibitem{chetelat97}
Ch\'etelat O., Merlet J-P., Myszkorowski P., and Longchamp R.
\newblock Globally convergent iterative algorithms for the coordinate
  transformations in the articulated mechanisms.
\newblock In {\em Syroco}, Nantes, \frenchmy{\Septembre}{1997}.

\bibitem{didrit98}
Didrit O., Petitot M., and Walter E.
\newblock Guaranteed solution of direct kinematic problems for general
  configurations of parallel manipulator.
\newblock {\em IEEE Trans. on Robotics and Automation}, 14(2):259--266,
  \frenchmy{\Avril}{1998}.

\bibitem{dietmaier98}
Dietmaier P.
\newblock The {Stewart-Gough} platform of general geometry can have 40 real
  postures.
\newblock In {\em ARK}, pages 7--16, Strobl,
  \frenchddate{29}{4}{\Juin}{\Juillet}{1998}.

\bibitem{faugere95}
Faug\`ere J.C. and Lazard D.
\newblock The combinatorial classes of parallel manipulators.
\newblock {\em Mechanism and Machine Theory}, 30(6):765--776,
  \frenchmy{\Aout}{1995}.

\bibitem{geng91}
Geng Z. and Haynes L.S.
\newblock Neural network for the forward kinematics problem of a {Stewart}
  platform.
\newblock In {\em IEEE Int. Conf. on Robotics and Automation}, pages
  2650--2655, Sacramento, \frenchdate{11-14}{\Avril}{1991}.

\bibitem{griffis89}
Griffis M. and Duffy J.
\newblock A forward displacement analysis of a class of {Stewart} platform.
\newblock {\em J. of Robotic Systems}, 6(6):703--720, 1989.

\bibitem{hunt78}
Hunt K.H.
\newblock {\em Kinematic geometry of mechanisms}.
\newblock Clarendon Press, Oxford, 1978.

\bibitem{hunt98}
Hunt K.H. and McAree P.R.
\newblock The octahedral manipulator: geometry and mobility.
\newblock {\em Int. J. of Robotics Research}, 17(8):868--885, 1998.

\bibitem{husty94-1}
Husty M.L.
\newblock An algorithm for solving the direct kinematic of
  {Stewart}-{Gough}-type platforms.
\newblock Research Report TR-CIM-94-7, Universit\'e McGill, Montr\'eal,
  \frenchdate{30}{\Juin}{1994}.

\bibitem{innocenti92-3}
Innocenti C. and Parenti-Castelli V.
\newblock Singularity-free evolution from one configuration to another in
  serial and fully-parallel manipulators.
\newblock In {\em 22nd Biennial Mechanisms Conf.}, pages 553--560, Scottsdale,
  \frenchdate{13-16}{\Septembre}{1992}.

\bibitem{lazard92}
Lazard D.
\newblock {Stewart} platform and {Gr\"obner} basis.
\newblock In {\em ARK}, pages 136--142, Ferrare,
  \frenchdate{7-9}{\Septembre}{1992}.

\bibitem{lazard94}
Lazard D. and Merlet J-P.
\newblock The (true) {Stewart} platform has 12 configurations.
\newblock In {\em IEEE Int. Conf. on Robotics and Automation}, pages
  2160--2165, San Diego, \frenchdate{8-13}{\Mai}{1994}.

\bibitem{merlet89-1}
Merlet J-P.
\newblock Manipulateurs parall\`eles, 4eme partie : mode d'assemblage et
  cin\'ematique directe sous forme polynomiale.
\newblock Research Report 1135, INRIA, \frenchmy{\Decembre}{1989}.


\bibitem{mourrain93}
Mourrain B.
\newblock The 40 generic positions of a parallel robot.
\newblock In Bronstein M., editor, {\em ISSAC'93}, ACM press, pages 173--182,
  Kiev (Ukraine), \frenchmy{\Juillet}{1993}.

\bibitem{nanua89}
Nanua P. and Waldron K.J.
\newblock Direct kinematic solution of a {Stewart} platform.
\newblock In {\em IEEE Int. Conf. on Robotics and Automation}, pages 431--437,
  Scottsdale, \frenchdate{14-19}{\Mai}{1989}.

\bibitem{raghavan91}
Raghavan M.
\newblock The {Stewart} platform of general geometry has 40 configurations.
\newblock In {\em ASME Design and Automation Conf.}, volume 32-2, pages
  397--402, Chicago, \frenchdate{22-25}{\Septembre}{1991}.

\bibitem{reboulet85}
Reboulet C. and Robert A.
\newblock Hybrid control of a manipulator with an active compliant wrist.
\newblock In {\em 3rd ISRR}, pages 76--80, Gouvieux, France,
  \frenchdate{7-11}{\Octobre}{1985}.

\bibitem{ronga92}
Ronga F. and Vust T.
\newblock {Stewart} platforms without computer?, 1992.
\newblock Preprint.

\bibitem{yee91}
Yee C.S. and Lim K.B.
\newblock Neural network for the forward kinematics problem in parallel
  manipulator.
\newblock In {\em IEEE Int. Joint Conf. on Neural Network}, pages 1699--1704,
  Singapore, \frenchdate{18-21}{\Novembre}{1991}.

\end{thebibliography}

\end{document} 
